\section{Assembly Wrapping: Mechanism and Integration}\label{sec:wrapping}
\subsection{Reading the Same Bytes from Different Starting Points}
Assembly wrapping relies on a property of x86 and x64: instructions have
variable lengths, and the starting byte determines how a decoder groups
the following bytes into instructions. The same stored bytes can therefore
have two interpretations when read from different starting points.

Consider two views of one region. In the \emph{outer view}, a decoder sees
an instruction followed by its immediate operand, a constant stored inside
that instruction. In the \emph{execution view}, part of that constant's
byte field is read as an instruction sequence. The established wrapping
technique uses this overlap to give the stored representation an outer
instruction structure while retaining an intended executable interpretation
within it~\cite{assembly_wrapping}.

Figure~\ref{fig:wrapping_views} shows these two readings of one byte field.
The term \emph{container} refers to the outer instruction; \emph{embedded
sequence} refers to the instructions represented inside its operand bytes.
These are two readings of shared bytes, rather than two copies of the
program. No decryption step is required to obtain the embedded interpretation.

\begin{figure}[t]
\centering
\begin{tikzpicture}[font=\small]
\node[anchor=east] at (0,0.5) {Stored bytes};
\draw[fill=gray!15] (0.15,0.15) rectangle (2.7,0.85);
\draw[fill=blue!8] (2.7,0.15) rectangle (7.4,0.85);
\node at (1.425,0.5) {Outer opcode};
\node at (5.05,0.5) {Immediate-operand field};
\node[anchor=east] at (0,1.65) {Outer view};
\node[draw,align=center,text width=6.9cm] at (3.775,1.65)
  {One instruction containing a constant operand};
\node[anchor=east] at (0,-0.9) {Execution view};
\node[draw,fill=blue!8,align=center,text width=4.1cm] (inside)
  at (5.05,-0.9) {Embedded instruction sequence};
\draw[-{Latex[length=2mm]},thick] (2.7,0.15) -- (2.7,-0.9) -- (inside.west);
\end{tikzpicture}
\caption{Conceptual views of a wrapped region. Starting at the outer
opcode interprets the operand as a constant; entering the embedded
sequence interprets some of those same bytes as instructions.
The drawing is not to scale and omits inter-region control transfers.}
\label{fig:wrapping_views}
\end{figure}

\subsection{What Zone Adds and What the Transformation Requires}
The contribution of this work is the integration and assessment of an
existing technique, with implementations for both architectures. Zone
starts from a prepared instruction stream, groups supported instructions
into wrapped regions, preserves the required execution relationships, and
updates the supported location-dependent references after layout. Some
instructions or local layouts cannot be placed in the main container
representation; the implementation retains these in separate, unwrapped
regions. An output may therefore contain both wrapped and ordinary code.

The x86 and x64 implementations share this purpose but use different
encodings and addressing rules. In particular, the available immediate
operand has a different width, affecting which instruction groups fit.
Both architectures are included in the functional evaluation.

Data preparation remains necessary. Wrapping does not determine whether
an apparent instruction is actually an embedded string or table. If that
earlier interpretation is wrong, construction may be refused or the output
may fail when executed. The following execution tests measure whether the resulting program
performs its expected task.

\subsection{How Wrapping Is Evaluated}
The evaluation follows three stages: construction of a binary, execution
of its functional test, and static scanning of the produced file.
Construction rate measures how many requested variants are available;
working rate measures how many perform the expected task; detection rate
describes the scanner responses. The full experiment uses all 102 original
payloads, with the same stages reported separately for x64 and x86.
